\chapter{Limitations and Future Work}\label{ch:limits}

\section*{Introduction}
\addcontentsline{toc}{section}{Introduction}

The previous chapter measured what this system does. This one is about what it
does not do, what that costs, and what would have to happen next. A fair few of
these limitations were not planned for --- I found them along the way, usually
by accident --- and a couple of them undercut things claimed earlier in this
report. That is exactly why they get their own chapter instead of a paragraph
buried at the end of the conclusion.

\section{Limitations}

\subsection{The Evaluation Is Narrow}

The golden set holds sixteen questions. That is enough to show that factual
retrieval and the safety behaviours hold up, and enough to show that
hallucination resistance does not. It is nowhere near enough to state an
accuracy rate with any real confidence, and any confidence interval in this
report would be meaningless.

There is also a more awkward problem: I wrote both the agent and the tests. No
matter how careful you are, the questions you think to ask are the failure modes
you already suspected. A genuinely independent set would almost certainly find
things this one does not.

And because the agent is not deterministic, one run measures one sample of its
behaviour rather than the behaviour itself. During this project the same
question was seen answering correctly once and incorrectly on a later attempt
with nothing changed in between. So a rate computed from a single pass is an
estimate with a variance nobody has measured.

\subsection{Grounding Filters, It Does Not Guarantee}

The verification layer catches two specific things: numbers that show up in an
answer but not in any query result, and claims that something is absent with no
zero-row query behind them. Both were built after watching the agent actually do
it, and both work.

Neither one touches the failure the evaluation found most often. Ask for a
metric the schema does not contain and the agent does not invent a number --- it
swaps in a neighbouring real one and reports it perfectly accurately. Call setup
success rate becomes something derived from dropped calls; customer lifetime
value becomes average revenue per user. Every number is traceable to a real
query, so every check passes happily. The system has no idea what the question
asked for, only what the answer contains, and that is a structural limit rather
than a missing rule I forgot to write.

\subsection{Synthetic Data Cuts Both Ways}

Synthetic data is the reason this project exists at all: it allowed a schema
modelled on a real industry standard with zero confidentiality exposure, and it
is the only reason objective ground truth for the evaluation is even possible.
On a live network the reference answer would itself be up for debate.

It also limits every result in here. Real operational data has missing values,
schema drift and duplicate identities that a generator does not reproduce unless
somebody sits down and writes them in. The clearest example is the churn work in
Chapter~\ref{ch:eval}: the heaviest term in the scorecard measures decline in
spend, and it separates absolutely nothing, because every simulated subscriber
follows the same revenue drift. There is no declining cohort in this world for a
decline term to find. A generator can produce data that looks perfectly
realistic and still contain none of the structure a model is supposed to detect
--- and there is no way to find that out except by going and looking, which is a
warning that applies to every result in this report.

\subsection{A Model Trained Elsewhere Was Deployed Without Checking}

The churn model was trained on a public dataset and pointed at this environment
without anybody verifying the two actually resembled each other. They do not:
eight of fifteen features diverge materially, one of them by a factor of three
hundred, and two are exact copies of two others. Its reported performance was
real and completely irrelevant.

This only came out because writing the evaluation chapter forced the question.
It is left here as a limitation rather than quietly patched, because the mistake
underneath it --- trusting a number produced on one distribution as evidence
about a different one --- is a lot more general than this project. And the
scorecard that replaced it is not validated either. It is defensible, it is
explainable, and it has never been measured against an outcome, because this
world does not produce any.

\subsection{One User, One Session}

Conversation state lives inside the server process. Two operators using the
system at the same time would share memory, and there is no authentication, no
per-user isolation and nothing survives a restart. Fine for a demonstration,
disqualifying for anything real.

\subsection{Things Nobody Measured}

Three things this report says nothing about, for the simple reason that nothing
was measured. \textbf{Latency} is asserted rather than benchmarked --- the fast
path and the semantic cache both exist to bring it down, and there is no number
anywhere saying by how much. \textbf{Cost} at operator scale was never
estimated; every single query is a cloud model invocation, and what that adds up
to across millions of subscribers is anybody's guess. \textbf{Security} of the
action pipeline was never analysed: the message broker carrying approved actions
has no authentication on it, and no threat model was ever written for a system
that can propose commercial actions against real customers.

\section{Future Work}

\subsection{A Serious Evaluation Harness}

Everything else on this list gets easier to justify once accuracy is a tracked
number instead of a feeling, which makes this the first thing to do. The harness
that exists is the skeleton. What it needs is scale --- somewhere around two
hundred questions covering simple lookups, multi-table joins, aggregations,
cross-domain questions and deliberately vague ones --- and independence, meaning
questions written by somebody who did not build the agent. Run it automatically
after every change and regression stops being invisible.

\subsection{Reining In the Non-Determinism}

Two things worth trying here, neither of them large. Asking the same question a
few times and comparing the query plans would surface the cases where the agent
quietly takes a different route, and disagreement between runs is itself a
useful hint that the question was vague to begin with --- which could feed
straight into the disambiguation gate that already exists. Separately, caching a
query plan once it has been verified would make repeat questions both faster and
stable.

\subsection{Grounding the Question, Not Just the Answer}

The substitution failure needs a check the current design simply cannot express:
comparing what was asked for against what came back. A workable version is
narrow enough --- pull the metric named in the question, verify it maps to a
real column or a documented derivation, and refuse when it does not. That would
have caught all three substitution failures, and it is really just the offer
catalogue fix from Chapter~\ref{ch:agent} generalised: hand the model the
authoritative vocabulary and check against it, instead of hoping.

\subsection{Actually Validating the Churn Scorecard}

The simulator already deactivates subscribers when they churn, which means it
has been producing labelled outcomes this whole time and nobody has used them.
Score the population, let a few simulated months go by, then measure which
flagged subscribers actually left. That would give the scorecard the validation
it currently does not have, and would test whether those weights are sensible or
just plausible. Cheapest path from an unvalidated heuristic to a measured one,
and it needs no external data at all.

\subsection{Towards Something Deployable}

Three things stand between this and something an operator could actually run.
Per-session state with authentication makes it multi-user. A schema-ingestion
step that reads an existing OSS/BSS catalogue and builds the knowledge graph
automatically removes the assumption that somebody hand-wrote the schema
knowledge in advance. And the action pipeline needs a proper audit trail: every
proposal, who approved or rejected it, and what actually got sent. None of that
is research. All of it is a prerequisite.

\subsection{A Richer World to Query}

Two gaps in the environment are worth closing, and both were found by work in
this report. The 2G layer carries no traffic at all right now and could model
the machine-to-machine bearer it actually is in real networks --- meters,
trackers, payment terminals --- as long as those devices are kept out of the
averages that describe human subscribers. And subscriber revenue could follow a
different trajectory per subscriber instead of one shared drift, so that a
declining cohort exists to be found. Without that, no churn signal based on
decline can be evaluated here at all.
